% ============================================================
% Section: Selection and justification of the evaluation dataset
% Requires: \usepackage{cite} or natbib/biblatex in the main
%           document preamble, and the references.bib file
%           (included at the end of this file) added to the
%           project bibliography.
% ============================================================

\section{Selection and Justification of the Evaluation Dataset}
\label{sec:dataset-selection}

In order to quantitatively evaluate the performance of the
fact-verification system proposed in this work, a reference dataset
(\textit{benchmark}) is required against which standard
classification metrics (precision, recall, and $F_1$) can be
measured relative to a verified \textit{ground truth}. This section
describes the process followed to select such a dataset, justifying
why the most widely cited benchmark in the automated fact-checking
literature --- FEVER \cite{thorne2018fever} --- was discarded in
favour of X-Fact \cite{gupta2021xfact}, and how this decision is
directly motivated by the architecture of the developed system.

\subsection{The Closed-Retrieval Paradigm: FEVER}
\label{subsec:fever}

FEVER (\textit{Fact Extraction and VERification}) \cite{thorne2018fever}
is, to date, the most influential dataset in automated
fact-verification research. It consists of 185,445 claims generated
through the controlled alteration of sentences extracted from
Wikipedia, each subsequently labelled by human annotators ---
without knowledge of the original sentence --- as
\texttt{SUPPORTED}, \texttt{REFUTED}, or \texttt{NOT ENOUGH INFO}.
For each supported or refuted claim, the sentence or set of
sentences from Wikipedia constituting the necessary evidence for
that verdict is additionally recorded.

The aspect of FEVER most relevant to this work is its
\textbf{closed-domain retrieval paradigm}: both the claim and its
correct evidence originate from a fixed, bounded corpus --- a static
dump of English Wikipedia --- and systems evaluated against this
benchmark typically index that same corpus to perform evidence
retrieval. The original paper itself characterises the task as the
combination of three subtasks evaluated jointly: document selection,
sentence selection within those documents, and final verdict
classification, all operating over the same closed corpus.

Direct multilingual extensions of FEVER also exist, such as XFEVER,
which automatically translate (via DeepL) the original claim-evidence
pairs into other languages, including Spanish. However, these
extensions inherit the same structural limitation as the original
dataset: the \textit{correct} evidence remains defined as a specific
sentence within a closed, translated corpus, rather than as any
valid evidence available on the open web.

\subsection{Why FEVER Is Not Suitable for This Work}
\label{subsec:discard-justification}

The system developed in this thesis --- described in detail in
Chapter~\ref{ch:methodology} --- does not perform verification
against a closed corpus. Its architecture is organised into six
sequential stages:

\begin{enumerate}
    \item \textbf{Admission}: ingestion and decomposition of the
    news article into atomic, verifiable claims.
    \item \textbf{Selection}: entity extraction and formulation of
    search queries from each claim.
    \item \textbf{Retrieval}: retrieval of candidate documents via
    SearXNG, a metasearch engine that queries multiple search
    engines over the open web in real time.
    \item \textbf{Ranking}: ordering of the retrieved sources by
    relevance to the claim.
    \item \textbf{LLM check}: reasoning over the ranked evidence by a
    language model (Llama 3.1, run locally via Ollama) to produce a
    preliminary verdict.
    \item \textbf{Recalibration}: adjustment of the verdict and its
    associated confidence, propagating the result to the article's
    overall verdict.
\end{enumerate}

This architecture --- referred to throughout this work as the
\textbf{proposed method} --- therefore performs
\textbf{open-domain retrieval}: no fixed corpus exists against which
claims are compared; instead, evidence is searched for \textit{on
the fly} on the live web at verification time, exactly as a human
fact-checker would.

Evaluating this system against FEVER or its multilingual variants
would introduce a serious methodological mismatch, for two reasons:

\begin{itemize}
    \item \textbf{Invalid-evidence bias.} The system could locate,
    through SearXNG, evidence that is perfectly valid and sufficient
    in a source other than the predefined Wikipedia sentence, and
    still be penalised for not matching the expected evidence
    verbatim. The resulting metric would systematically underestimate
    the system's true performance.
    \item \textbf{Partial evaluation of the architecture.} Since
    FEVER fixes the source corpus in advance, evaluation against this
    benchmark would measure almost exclusively the \textit{LLM check}
    stage, leaving the \textit{Selection}, \textit{Retrieval}, and
    \textit{Ranking} stages unevaluated --- precisely the stages that
    constitute the most original, and also the most fragile,
    contribution of the system, as discussed in
    Chapter~\ref{ch:results}.
\end{itemize}

For these reasons, FEVER/XFEVER was discarded as the primary
evaluation benchmark: not out of convenience or expected performance,
but because the dataset's own design assumes a (closed) retrieval
paradigm incompatible with the (open) retrieval paradigm underlying
the proposed system.

\subsection{X-Fact as an Open-Retrieval Benchmark}
\label{subsec:xfact}

As an alternative, \textbf{X-Fact} \cite{gupta2021xfact} was
selected: the largest publicly available multilingual dataset for
the verification of naturally occurring, real-world claims. X-Fact
contains 31,189 claims across 25 languages --- including Spanish ---
drawn from professional fact-checking organisations and labelled by
expert fact-checkers using a fine-grained veracity scale (not limited
to FEVER's supported/refuted/not-enough-info triad).

The decisive structural difference with respect to FEVER is that
X-Fact's claims are \textbf{real-world claims} made in public
discourse (statements by public figures, news headlines, viral
posts), not synthetic sentences derived from Wikipedia, and their
verification --- both by the original human fact-checkers and by the
baseline models proposed by the authors --- relies on evidence
retrieved via search engines over the open web, not against a
predefined closed corpus. The authors themselves report that their
best baseline model, which combines the claim text with metadata and
evidence retrieved through a search engine, achieves an $F_1$ of
approximately 40\%, positioning X-Fact as a challenging benchmark
even for systems that already operate under the same open-retrieval
paradigm as the method proposed in this work.

This alignment of paradigms --- open retrieval in the evaluation
dataset and open retrieval in the evaluated system --- is what makes
X-Fact a methodologically coherent choice: it allows the full
pipeline (Selection, Retrieval, Ranking, and LLM check jointly) to be
measured, not only its final reasoning stage, and situates the
results obtained within a framework comparable to the state of the
art reported by Gupta and Srikumar \cite{gupta2021xfact} for the
Spanish-language subset.

\subsection{Domain Limitation and a Custom Validation Set}
\label{subsec:custom-gold-set}

Nevertheless, X-Fact does not fully capture the nature of this
work's application domain either. Its claims originate predominantly
from political discourse and controversial public statements, while
the developed system is applied to editorial articles from a
positive-news outlet, with claims combining verifiable facts
(figures, dates, events) with analytical or cultural
interpretations --- a type of claim virtually absent from X-Fact and
which, as documented in Chapter~\ref{ch:results}, exposes a real
limitation of the system in its current version.

For this reason, the quantitative evaluation in this work relies on
\textbf{two complementary sources}: X-Fact \cite{gupta2021xfact}, as
a standardised benchmark enabling comparison of the proposed method
against the state of the art in multilingual, open-retrieval fact
verification; and a custom validation set, built from claims
extracted from real articles published by the outlet under study and
manually labelled, specifically designed to expose and measure the
system's behaviour on interpretive claims --- the real use case for
which the system was conceived, and one that no existing academic
benchmark adequately covers.

% ============================================================
% Corresponding BibTeX entries (add to references.bib)
% ============================================================
% @inproceedings{thorne2018fever,
%   title     = {{FEVER}: a Large-scale Dataset for Fact Extraction and {VER}ification},
%   author    = {Thorne, James and Vlachos, Andreas and Christodoulopoulos, Christos and Mittal, Arpit},
%   booktitle = {Proceedings of the 2018 Conference of the North American Chapter of the Association for Computational Linguistics: Human Language Technologies, Volume 1 (Long Papers)},
%   pages     = {809--819},
%   year      = {2018},
%   address   = {New Orleans, Louisiana},
%   publisher = {Association for Computational Linguistics},
%   doi       = {10.18653/v1/N18-1074},
%   note      = {arXiv:1803.05355},
%   url       = {https://arxiv.org/abs/1803.05355}
% }
%
% @inproceedings{gupta2021xfact,
%   title     = {{X-Fact}: A New Benchmark Dataset for Multilingual Fact Checking},
%   author    = {Gupta, Ashim and Srikumar, Vivek},
%   booktitle = {Proceedings of the 59th Annual Meeting of the Association for Computational Linguistics and the 11th International Joint Conference on Natural Language Processing (Volume 2: Short Papers)},
%   pages     = {675--682},
%   year      = {2021},
%   address   = {Online},
%   publisher = {Association for Computational Linguistics},
%   doi       = {10.18653/v1/2021.acl-short.86},
%   url       = {https://aclanthology.org/2021.acl-short.86}
% }